\documentclass[sigplan,10pt]{acmart}

\setcopyright{none}
\renewcommand\footnotetextcopyrightpermission[1]{}
\settopmatter{printfolios=true}
\pagestyle{plain}

\acmConference[ATC '26]{2026 ACM SIGOPS Annual Technical Conference}{November 15--18, 2026}{Hyatt Hotel, Shatin, Hong Kong}
\acmYear{2026}
\acmISBN{}
\acmDOI{}

\title[ActPlane]{ActPlane: OS-Enforced Agent Harnesses}

\author{Paper \#XXX}
\affiliation{%
  \institution{Anonymous Submission}
  \city{}
  \country{}}

\begin{abstract}
Modern coding agents are governed largely by natural-language instructions such
as \texttt{CLAUDE.md} and \texttt{AGENTS.md}. These instructions express real
workflow and data-handling constraints, but they are probabilistic: an otherwise
cooperative agent may forget them, invoke a shell, spawn a subprocess, or use an
SDK path that bypasses a tool-level guard. Existing agent guardrails therefore
sit at the wrong layer for constraints that must hold across the agent's whole
execution.

ActPlane is an OS-level harness for AI agents. It compiles declarative
information-flow policies into an eBPF/BPF-LSM enforcement backend that tracks
typed labels across process, file, and endpoint objects; checks source-to-sink
rules at kernel hooks; and returns corrective feedback to the agent so that a
failed action can be repaired rather than blindly retried. The system targets
cooperative-but-forgetful agents, not only adversarial code. Its semantics
therefore distinguish LSM-only pre-operation \emph{block}, tracepoint-compatible
\emph{audit}, and harness-level \emph{kill}. This paper describes ActPlane's
policy language, kernel design, feedback bridge, and an evaluation plan that
will measure bypass coverage, agent recovery, false positives, overhead, and
real-policy coverage. Evaluation result cells are intentionally left as TBDs
until the experiments are run.
\end{abstract}

\begin{CCSXML}
<ccs2012>
 <concept>
  <concept_id>10010520.10010553.10010562</concept_id>
  <concept_desc>Computer systems organization~Embedded and cyber-physical systems</concept_desc>
  <concept_significance>300</concept_significance>
 </concept>
 <concept>
  <concept_id>10002978.10003022.10003023</concept_id>
  <concept_desc>Security and privacy~Software and application security</concept_desc>
  <concept_significance>300</concept_significance>
 </concept>
 <concept>
  <concept_id>10011007.10010940.10010941.10010949</concept_id>
  <concept_desc>Software and its engineering~Software safety</concept_desc>
  <concept_significance>300</concept_significance>
 </concept>
</ccs2012>
\end{CCSXML}

\ccsdesc[300]{Computer systems organization~Embedded and cyber-physical systems}
\ccsdesc[300]{Security and privacy~Software and application security}
\ccsdesc[300]{Software and its engineering~Software safety}

\keywords{AI agents, eBPF, BPF-LSM, information-flow control, provenance, taint tracking, corrective feedback}

\begin{document}
\maketitle

\section{Introduction}

AI coding agents increasingly operate as long-running processes with access to
shells, source trees, package managers, and external APIs. Project maintainers
steer these agents with instruction files such as \texttt{CLAUDE.md} and
\texttt{AGENTS.md}: do not push directly to \texttt{main}; run tests before
committing; never expose secrets; use only a sanctioned review tool before
destructive operations. These are not style preferences. They are behavioral
constraints that must hold across the agent's entire execution.

A cooperative agent will generally follow these constraints---but it forgets.
In long contexts, attention is diluted, earlier instructions drop out of the
effective window, and the agent takes a shortcut that violates a rule it
previously acknowledged. The agent is not adversarial; it is
\emph{cooperative but forgetful}. It would benefit from a mechanism that
makes its own declared constraints durable---surviving context loss, tool
switching, and subprocess indirection---and that reminds it when a constraint
is about to be violated.

\subsection{The Intent--Behavior Gap}

We distinguish three levels at which an agent's execution can be described,
following AgentSight's observation of a ``semantic gap'' between agent intent
and low-level actions~\cite{agentsight}.

\textbf{Intent.} What the agent declares it will and will not do: goals
(``fix this bug'') and self-constraints (``I must not leak secrets'').
Intent is expressed in natural language and, in ActPlane's design, formalized
in a compact DSL that the agent itself can write and maintain.

\textbf{Action.} The tool calls the agent issues through its runtime:
\texttt{read\_file}, \texttt{run\_command}, \texttt{edit\_file}. This is the
interface between the agent and the outside world, mediated by the agent
framework.

\textbf{Behavior.} The actual OS-level operations that result: \texttt{open},
\texttt{execve}, \texttt{connect}, \texttt{fork}. Behavior is where effects
really happen---files are read, processes are spawned, data reaches the
network.

The core problem is that the mapping from action to behavior is
\emph{many-to-many and opaque}. A single tool call
\texttt{run\_command("make")} can produce hundreds of system calls. The same
OS operation---\texttt{execve("git", ["push"])}---can be reached through a
direct tool call, a \texttt{bash~-c} string, a Python subprocess, or a
compiled helper binary. An action-level guard sees only the tool call; it
cannot predict or intercept the behavior it produces.

\subsection{Why Single-Level Enforcement Is Insufficient}

Enforcement can be placed at any of the three levels. Each alone is
insufficient.

\textbf{Intent-level} enforcement (prompt instructions, \texttt{CLAUDE.md})
is probabilistic: the agent may forget the constraint or override it in
reasoning.

\textbf{Action-level} enforcement (tool-call guards such as AgentSpec and
Progent~\cite{wang2025agentspec,shi2025progent}) is deterministic but
bypassable: it checks tool calls, not the OS behavior they produce. When the
agent routes an effect through a subprocess, an SDK, or a direct system
interface, the guard sees nothing.

\textbf{Behavior-level} enforcement (syscall filters, sandboxes, eBPF
enforcers) is unbypassable---every effect crosses the kernel boundary. But
isolated behavior-level enforcement returns only \texttt{-EPERM} or
\texttt{SIGKILL}; the agent receives no explanation and cannot recover.

The gap between intent and behavior is not bridged by any single level.
Intent-level constraints are not connected to behavior-level enforcement;
behavior-level enforcement is not connected back to intent-level feedback.

\subsection{ActPlane: Bridging Intent and Behavior}

ActPlane closes this gap in both directions.

\textbf{Downward: intent to behavior.} The agent declares behavioral
contracts in a DSL---voluntary confinement, analogous to \texttt{pledge()} in
OpenBSD, where a program restricts its own capabilities because it knows it
may be exploited. A coding agent can write and maintain this DSL just as it
edits any other project file. ActPlane compiles these declarations into
labeled information-flow rules and loads them into an eBPF/BPF-LSM backend.
Named labels propagate across process, file, and network objects at kernel
hooks, tracking which sources have influenced which subjects.

\textbf{Horizontal: across objects.} Labels encode execution history at the
behavior level. When a process reads a file labeled \texttt{SECRET}, the
process acquires that label. When it forks, the child inherits the label.
This cross-object propagation turns individual OS operations into checkable
behavioral state---enabling constraints like ``a process that has read
\texttt{.env} must not connect to external endpoints.''

\textbf{Upward: behavior to intent.} When a violation is detected, the
kernel reports the rule, the target, and the effect. ActPlane translates
this into corrective feedback delivered to the agent's context: which
constraint was violated, why the current operation conflicts with it, and
what alternative path (running a redactor, going through a review tool,
executing tests first) would satisfy the contract. The agent receives not
\texttt{Permission denied} but a reminder of its own declared intent.

The term \emph{agent harness} has come to mean everything around the
model---tool dispatch, memory, orchestration, guardrails, and
feedback~\cite{langchain-harness,fowler-harness}. Current harness
implementations place all of these components, including enforcement and
feedback, at the application layer. ActPlane extends the harness to the OS
level for the components that cannot work at the application layer alone:
behavioral enforcement and corrective feedback over the agent's actual
OS-level behavior. It is not a replacement for the application-layer harness
but a complement---the part that must reach below the tool API.

This design is not a new enforcement mechanism---CamQuery showed that
cross-channel label propagation and enforcement can be done in the
kernel~\cite{pasquier2018camquery}. Nor is corrective feedback for agents
new---AgentSpec demonstrated it at the tool-call
layer~\cite{wang2025agentspec}. And letting agents participate in their own
governance through information-flow control is already practiced at the
application layer by systems such as FIDES and
CaMeL~\cite{costa2025fides,debenedetti2025camel}. ActPlane's contribution is
bringing these capabilities together \emph{at the OS level}: cross-channel
labeled information-flow enforcement on the modern eBPF/BPF-LSM substrate,
driven by agent-declared behavioral contracts, with a feedback path that
connects kernel-detected violations back to the agent's intent---where
constraints survive context loss and cannot be bypassed below the tool API.

\subsection{Contributions}

This paper makes three contributions.

\begin{enumerate}
\item It defines the \emph{intent--behavior gap} for AI agents---the
disconnect between behavioral constraints declared at the intent level and
their enforcement at the OS level---and grounds it with a corpus study of
instruction files from 50 popular projects, showing that OS-enforceable
behavioral constraints are pervasive.

\item It presents ActPlane, a system that compiles agent-declared behavioral
contracts into eBPF/BPF-LSM labeled information-flow enforcement. The system
propagates labels across process, file, and network objects, checks
source-to-sink rules at kernel hooks, and returns corrective feedback that
reconnects behavior-level violations to intent-level remediation.

\item It evaluates ActPlane on bypass coverage across five tool paths, agent
recovery rate under four feedback conditions, false positives under benign
workloads, and per-event overhead, comparing against action-level and
behavior-level baselines.
\end{enumerate}

\section{Motivation and Requirements}

\subsection{Natural-Language Harnesses Are Real}

We built a snapshot and analysis pipeline for
\texttt{CLAUDE.md} and \texttt{AGENTS.md} files from popular software projects.
The preliminary extraction shows that these files contain many imperative rules,
but only a subset are OS-observable behavioral constraints. This distinction is
important. A rule such as ``never use \texttt{any}'' is a coding-style rule and
is outside ActPlane's scope. A rule such as ``never push directly to
\texttt{main}'' or ``run tests before committing'' has an observable process
and syscall footprint and can be enforced as a harness contract.

The paper will use the corpus study for two purposes. First, it grounds the
motivation: developers already write behavioral constraints for agents. Second,
it derives evaluation workloads: only rules that are observable, expressible,
robust, and temptable are converted into policy/task pairs. The final paper
should report both the unbiased corpus denominator and the filtered evaluation
subset, rather than selecting only examples that ActPlane can handle.

\subsection{Why Tool-Layer Guards Are Insufficient}

Agent frameworks and permission systems can intercept tool calls and apply
deterministic policies~\cite{wang2025agentspec,shi2025progent}. This is useful,
but it assumes that effects pass through the guarded tool boundary. A coding
agent violates that assumption whenever it invokes a shell, starts a subprocess,
loads a library, or talks to an external service from generated code. The same
semantic action, for example invoking \texttt{git push --force}, can appear as a
direct shell command, a nested \texttt{bash -c} string, a Python subprocess, or
a helper binary. A tool-layer guard sees only the outer tool invocation.

The OS sees the effect itself. Every exec, file open, unlink, rename, truncate,
and socket connect crosses a kernel boundary. If a policy is phrased over those
effects and over the lineage of the agent process tree, then the policy applies
regardless of the high-level route the agent selected.

\subsection{Requirements}

ActPlane is designed around five requirements.

\textbf{R1: Whole-execution scope.} A rule must apply to the agent's
descendants, not just to the top-level process or the initial tool call.

\textbf{R2: Cross-object provenance.} Many rules are about information flow,
not a single event. Reading a secret file should taint the reader; writing a
derived file should taint that file; reading the derived file later should taint
the next process.

\textbf{R3: Effect semantics that match the backend.} \texttt{block} means
pre-operation denial and is implemented only by BPF-LSM hooks. Tracepoints do
not support \texttt{block}; they can report \texttt{audit} rules or terminate
\texttt{kill} rules.

\textbf{R4: Corrective feedback.} The harness should not only fail a violating
action. It should tell the agent which rule fired, why, and what alternative or
precondition can allow progress.

\textbf{R5: Deployable mechanism.} The implementation should use upstream Linux
interfaces: eBPF for dynamic loading and BPF-LSM for pre-operation deny, rather
than a custom kernel module.

\section{Policy Model}

ActPlane policies are object-level information-flow rules over three node
types: processes, files, and endpoints. A process node represents a Linux task
identified by pid and command identity. A file node is currently keyed by path.
An endpoint node is currently keyed by IPv4 address. The model is intentionally
coarser than byte-level dynamic taint tracking~\cite{kemerlis2012libdft}; it
trades precision for low-overhead whole-system coverage.

\subsection{Labels and Propagation}

Each node carries a bitset of labels. A source declaration introduces labels:
an exec source labels processes that execute a matching binary, a file source
labels matching files, and an endpoint source labels matching network
endpoints. Labels propagate through fixed transfer functions:

\begin{itemize}
\item fork: a child process inherits the parent's labels;
\item exec: a process absorbs labels from the executed file and matching exec
sources;
\item read: a process absorbs labels from the file it reads;
\item write: a file absorbs labels from the process writing it;
\item connect: an endpoint records labels from the connecting process.
\end{itemize}

This is the provenance invariant ActPlane maintains: if a process carries label
\texttt{SECRET}, then information from some \texttt{SECRET} source has reached
that process through observed OS-level flow edges.

\subsection{Rules, Gates, and Effects}

A rule contains one or more deny clauses. A clause specifies an operation
(\texttt{exec}, \texttt{read}, \texttt{write}, \texttt{unlink},
\texttt{connect}, or \texttt{open}), a target pattern, a Boolean expression over
labels, and an optional condition such as a target allow-list, a lineage gate,
or an after gate.

\begin{figure}[t]
\begin{verbatim}
label AGENT
source SECRET = file "**/.env"

rule no-secret-egress:
  deny connect endpoint "*" if SECRET
  effect block
  reason "secret-derived data must stay local"
  remediation "run the redactor or ask the user"

rule no-agent-git:
  deny exec "**/git" if AGENT
  effect kill
  reason "agent runs must not invoke git"
\end{verbatim}
\caption{A shortened ActPlane policy. The first rule asks for LSM-only
pre-operation deny. The second rule asks for harness-level termination and can
fire from tracepoint observation.}
\Description{A code listing showing two ActPlane rules: one blocking secret
egress and one killing git execution by an agent.}
\label{fig:policy-example}
\end{figure}

Rules have explicit effects.

\textbf{Audit.} The kernel reports the violation and allows the action to
proceed. This is useful for soft workflow guidance and for measurement.

\textbf{Block.} A BPF-LSM hook returns \texttt{-EPERM} before the operation is
committed. \texttt{block} is not supported in tracepoint mode and is not
silently downgraded to \texttt{audit} or \texttt{kill}.

\textbf{Kill.} The kernel sends \texttt{SIGKILL} to the current task. If a
pre-operation LSM hook is active, ActPlane can also return \texttt{-EPERM};
from a tracepoint, \texttt{kill} is harness-level termination rather than
pre-operation denial.

If several supported rules match the same event, ActPlane chooses the strongest
effect in the order \texttt{kill > block > audit}. Unsupported effects are
ignored by that backend, so a tracepoint \texttt{audit} rule can still fire even
when an unsupported \texttt{block} rule also matches.

\subsection{Why This Is a Harness DSL}

The language is deliberately not a general-purpose MAC policy. Its primitives
map to agent workflows: ``after tests,'' ``through a review tool,'' ``data
derived from untrusted input,'' and ``read-only sub-agent.'' The grammar is
small because policy authors should be able to translate natural-language
instructions into a few recurring source/sink/gate shapes.

\section{Design}
\label{sec:design}

\sys{} is an OS-level harness for agent policy compliance. Its central
design choice is to let agents and project owners express policy in terms
of task and repository intent, while enforcing the resulting constraints
below the tool layer. This separates the context needed to instantiate a
policy from the mechanism that observes system-level behavior.

The design challenge is to keep these two sides connected without
collapsing one into the other: a policy must be concrete enough for
deterministic OS-level enforcement, but expressive enough to encode
project and task intent supplied by the agent. It must also preserve
cross-event state and prevent lower-authority policy from weakening
higher-authority constraints.

\subsection{Threat Model and Assumptions}

\sys{} targets a cooperative but unreliable agent. The agent is expected
to declare task intent and project constraints honestly, but may fail to
follow them because of planning errors, prompt injection, shell-script
side effects, or opaque tool behavior. The system does not target a
malicious process that attacks the kernel, exploits the loader, or
deliberately evades policy through adversarial obfuscation.

The trusted computing base consists of the compiler, the eBPF enforcement
engine, the runner, and higher-authority policies. The agent's process
tree, including tools, shells, subprocesses, and sub-agents, is treated as
the monitored domain. Out of scope are kernel compromise, root or
\texttt{CAP\_BPF} compromise, deliberate side channels, and semantic
content checks that are not observable at process, file, or network
edges.

\subsection{System Overview}

The empirical study in \S\ref{sec:empirical} gives three design
requirements. First, policy must be \emph{agent-facing}: many directives
refer to project- or task-specific concepts that only the agent can
resolve in context. Second, enforcement must be \emph{system-level}:
directives often constrain commands, files, subprocesses, and network
connections rather than final model text. Third, enforcement must be
\emph{stateful}: many directives depend on previous actions, such as
whether tests ran after a source change or whether data from one source
later reached another sink.

\sys{} implements these requirements with three abstractions. A compact
DSL lets agents and maintainers bind natural-language directives to
source, sink, and condition patterns. A labeled information-flow model
maintains cross-event state across processes, files, and network
endpoints. A layered authority model lets lower-authority policy add or
narrow constraints without weakening platform or project invariants.

\subsection{Policy DSL}

The DSL is the boundary between intent-level instructions and
system-level enforcement. A policy contains sources that introduce labels,
rules that match labeled events at sinks, and optional gates that express
ordering or controlled declassification. This keeps the agent-facing
surface small: translating an instruction usually requires producing one
declarative rule rather than constructing the backend representation used
by the kernel.

Sources bind labels to system objects. An \texttt{exec} source labels
processes that execute a matching binary, a \texttt{file} source labels
matching paths, and a network source labels matching endpoints. Rules bind
an effect such as \texttt{notify}, \texttt{block}, or \texttt{kill} to an
operation such as \texttt{exec}, \texttt{read}, \texttt{write},
\texttt{unlink}, or \texttt{connect}. Conditions express common agent
workflow constraints: a target allow-list, a lineage gate, or a temporal
\texttt{after ... since ...} gate.

\begin{figure}[t]
\begin{verbatim}
source AGENT = exec "claude"
source SECRET = file "**/.env"

# Per-event: from CoplayDev/unity-mcp
rule no-push-main:
  kill exec "git" "push" "main" if AGENT
  because "Do not push to main; open a PR instead."

# Cross-event: from chenhg5/cc-connect
rule tests-before-commit:
  block exec "git" "commit"
    if AGENT unless after exec "go" since write "**/*.go"
  because "Run `go test ./...` before committing."
\end{verbatim}
\caption{Two \sys{} rules drawn from real projects. The first is a
per-event rule; the second is a cross-event rule that requires temporal
state.}
\Description{Code listing of two \sys{} rules from real projects.}
\label{fig:policy-example}
\end{figure}

Figure~\ref{fig:policy-example} shows the two main rule shapes. The first
case is local to one event: if the agent process tries to push to
\texttt{main}, the event matches immediately. The second case depends on
history: committing is compliant only if the required test command ran
after the relevant source write. This distinction is why the DSL exposes
both per-event predicates and cross-event gates.

The DSL is intentionally narrower than a general mandatory access-control
language. It focuses on policy shapes that recur in agent instruction
files: do not use a command, do not write outside a scope, do not connect
to a sink after reading a source, run a validation command before a
release action, and route sensitive data through an approved tool. The
compiler lowers these rules into fixed-size kernel tables; policy authors
do not need to reason about BPF maps, verifier constraints, or binary
layout.

\subsection{Labeled Information-Flow Model}

\sys{} represents policy state as labels on OS objects. Processes, files,
and network endpoints each carry a 64-bit label mask. A source declaration
adds labels when an object matches the source pattern. Let $L(x)$ denote
the label mask of object $x$. Labels propagate through observed
system-level edges:

\begin{itemize}
\item \textbf{fork}$(p \to c)$: $L(c) := L(c) \cup L(p)$.
\item \textbf{exec}$(p, f)$: $L(p) := L(p) \cup L(f) \cup S(f)$, where
$S(f)$ is the set of source labels matching $f$.
\item \textbf{read}$(p, f)$: $L(p) := L(p) \cup L(f)$.
\item \textbf{write}$(p, f)$: $L(f) := L(f) \cup L(p)$.
\item \textbf{connect}$(p, e)$: $L(e) := L(e) \cup L(p)$.
\end{itemize}

These transfer functions maintain the key invariant: if an object carries
label $\ell$, then information from a source tagged with $\ell$ has
reached that object through observed process, file, or network edges.
Rules then check whether an event's subject or target contains required
labels, contains forbidden labels, or satisfies a temporal gate.

The model operates at object granularity rather than byte granularity,
trading precision for low-overhead whole-session coverage. This is a
deliberate fit for agent harnesses. Repository instructions usually refer
to objects such as commands, files, directories, repositories, endpoints,
and workflow steps, not individual bytes. Object-level labels can express
cross-event compliance patterns that tool-call filters miss, while
avoiding the complexity of dynamic byte-level taint
tracking~\cite{kemerlis2012libdft}.

Labels are monotonic by default: propagation adds labels and does not
remove them. When a project requires controlled release, the DSL can name
an explicit gate, such as a redaction command or review tool, that clears
specified labels for subsequent events. This makes declassification a
declared policy action rather than an implicit side effect of the agent's
plan.

\subsection{Layered Policy Authority}

Agent-authored policy is useful only if it cannot weaken stronger
constraints. \sys{} therefore treats policy as a layered authority stack:
platform/operator policy, project-maintainer policy, and agent/session
policy. Lower layers may add rules, specialize targets, add labels, or
narrow scope. They may not delete, rewrite, declassify, or weaken rules
introduced by a higher-authority layer.

\sys{} separates a pure IFC rule from the runtime binding that activates
it. A rule defines a condition, an effect, and a reason. A domain is the
policy boundary for a process tree; each process belongs to one domain,
and each event is checked against that domain's effective policy. Files
and endpoints are IFC objects that carry labels, not policy domains.

Bindings attach rules to domains as either \emph{locked} or
\emph{default}. Locked bindings are inherited security invariants and
cannot be disabled by child domains. Default bindings are templates that a
child domain may keep, disable for itself, or replace with stricter local
policy. This gives the agent room to adapt policy to a task while
preserving the authority of platform and project constraints.

The prototype resolves the active authority stack when policy is compiled
or loaded. The kernel receives the resulting rule tables and evaluates
events against them; it does not parse YAML, resolve project context, or
interpret authority metadata on the event path. This preserves the design
separation: agents and maintainers supply context and policy structure,
while the OS-level engine enforces the compiled result consistently across
tools and subprocesses.

\section{Implementation}
\label{sec:implementation}

\sys{} is implemented as a Rust compiler/runner and an eBPF enforcement
engine. The compiler parses the DSL, allocates label bits, lowers Boolean
label expressions and path/network patterns, and emits a fixed-size
C-compatible configuration consumed by the eBPF loader. Human-readable
rule names and reasons remain in userspace metadata; the kernel receives
compact rule tables and emits rule identifiers when events match.

The eBPF engine attaches to process, file, and network hooks. It maintains
label state for processes, files, and endpoints in BPF maps, applies the
transfer functions above, and checks the compiled rules on each observed
event. BPF-LSM hooks provide pre-operation enforcement for blocking
actions; tracepoints provide observation and post-operation termination
where LSM support is unavailable. Userspace does not re-detect policy
violations: it loads policy, starts the target process, and formats
kernel-reported matches.

The runner exposes this mechanism to agents through
\texttt{actplane run -- <cmd>}. It discovers the project policy file,
compiles and loads the policy before the target starts, seeds the target
process tree with the agent label, and records kernel matches in a
feedback file that compatible agent hooks can relay into the next turn.
This plumbing is intentionally thin: the policy decision remains in the
kernel, while the runner only connects project policy, target execution,
and agent-facing diagnostics.

\section{Implementation Status}

ActPlane is implemented as a Rust collector and an eBPF loader/program pair.
The current codebase builds a single \texttt{actplane} binary that embeds the
compiled loader, extracts it at runtime, compiles project policy YAML, and runs
a target command under the policy harness.

\subsection{Runner}

The command \texttt{actplane run -- <cmd>} discovers \texttt{actplane.yaml}
upward from the current directory unless \texttt{--policy} or \texttt{--rule} is
specified. The policy YAML contains a \texttt{policy: |} DSL block and an
optional feedback path. Run mode requires an \texttt{AGENT} label; before the
target begins executing, the loader seeds the stopped target pid with that label
so all descendants inherit the agent identity.

The runner prepares \texttt{.actplane/last-violation.txt} and a hook state file,
exports their paths to the child, waits until the loader reports readiness, then
continues the target process. This ensures the first target action is observed.

\subsection{BPF Backends}

The eBPF program attaches to process lifecycle tracepoints for fork, exec, and
exit, to syscall tracepoints for observation-compatible file and connect events,
and to BPF-LSM hooks for pre-operation enforcement. The loader checks whether
the running kernel has the \texttt{bpf} LSM active. If not, it uses tracepoint
mode. In tracepoint mode, supported effects are exactly \texttt{audit} and
\texttt{kill}; \texttt{block} is unsupported because there is no pre-operation
verdict to return.

\subsection{Current Limitations}

The implementation deliberately exposes several limitations rather than hiding
them behind ambiguous semantics.

\textbf{Path identity.} File labels are currently keyed by path hash. This is
sufficient for the prototype but should move toward inode/device identity for a
stronger claim.

\textbf{Endpoint identity.} Kernel connect matching is numeric IPv4. Hostname,
DNS, SNI, and HTTP-layer policies require userspace assist or additional
instrumentation.

\textbf{Exec arguments.} \texttt{@arg} predicates are currently observed after
exec unless an argv cache is added before \texttt{bprm\_check\_security}. Thus
\texttt{effect block} with an argv predicate is not available through the
tracepoint path; \texttt{effect audit} can report and \texttt{effect kill} can
terminate after exec.

\textbf{Precision.} Labels propagate at object granularity. This can over-taint
compared with byte-level DIFT systems. The evaluation must therefore measure
false positives and validate declassification/gate paths, not only violation
coverage.

\subsection{Engineering Invariants}

The codebase now follows three invariants that matter for the paper:

\begin{itemize}
\item The only kernel output for policy violations is
\texttt{TAINT\_VIOLATION}; ordinary tracing noise is not part of the harness
interface.
\item Each rule explicitly declares its effect. There is no global fallback that
converts \texttt{block} into \texttt{kill}.
\item The post-tool feedback adapter is a transport bridge, not a second policy
engine.
\end{itemize}

\section{Evaluation Plan}
\label{sec:evaluation}

The evaluation is organized around the following reproducible experiments.
Result numbers are not filled in until the corresponding experiment is run with
fixed agent versions, logs, and statistical analysis.

\subsection{Research Questions}

The evaluation targets six questions.

\begin{description}
\item[RQ1: Bypass coverage.] Does OS-level enforcement catch the same policy
violation through direct tool calls, shell strings, subprocesses, SDK paths, and
direct syscalls?
\item[RQ2: Corrective feedback.] Does semantic feedback improve task completion
and reduce repeated violations compared with failure without a reason?
\item[RQ3: False positives.] Do benign tasks and sanctioned gate/declassify
paths proceed without unnecessary failure?
\item[RQ4: Overhead.] What are the per-event and end-to-end costs of label
propagation and rule checks?
\item[RQ5: Real-policy coverage.] How many real instructions from the corpus
are observable, expressible, robust, and temptable?
\item[RQ6: Baselines.] Which capabilities are unique relative to tool-layer
agent guardrails and kernel provenance/enforcement systems?
\end{description}

\subsection{Experiments}

\begin{table*}[t]
\centering
\caption{Planned evaluation matrix. Result fields remain \textsc{TBD}.}
\Description{A table listing planned experiments and planned result fields.}
\begin{tabular}{p{0.12\textwidth}p{0.26\textwidth}p{0.26\textwidth}p{0.24\textwidth}}
\toprule
Experiment & Method & Primary metrics & Result status \\
\midrule
E1: Cross-path coverage &
Run the same forbidden action through tool call, nested shell, subprocess, SDK/helper, and direct syscall triggers. Compare against tool-layer baselines. &
Coverage matrix over policy by path; feedback delivery. &
\textsc{TBD} \\
\midrule
E2: Feedback loop &
Run Codex and Claude-style agents under prompt-only, audit, fail-without-reason, and fail-with-feedback conditions. &
Task completion, repeated violations, repair iterations, token/turn cost. &
\textsc{TBD} \\
\midrule
E3: False positives &
Run allowed tasks and sanctioned paths: redaction, test-before-commit, confirmation, and lineage mediation. &
False positive rate, allowed-path success. &
\textsc{TBD} \\
\midrule
E4: Overhead &
Microbench fork/exec/read/write/connect and macrobench real agent tasks with and without ActPlane. &
p50/p99 latency, wall time, memory/map usage, rule-table scaling. &
\textsc{TBD} \\
\midrule
E5: Corpus coverage &
Double-code real instruction-file rules with D1--D7, report agreement, and derive policies/tasks mechanically. &
Observable, expressible, robust, temptable fractions; Cohen's kappa. &
\textsc{TBD} \\
\midrule
E6: Baseline comparison &
Compare against CamQuery, Tetragon, Progent, AgentSpec, and built-in agent sandboxes where runnable. &
Layer, enforce/detect, cross-channel taint, bypass resistance, feedback. &
\textsc{TBD} \\
\bottomrule
\end{tabular}
\label{tab:eval-plan}
\end{table*}

\subsection{Feedback Conditions}

The agent-facing experiment uses four conditions to isolate the contribution of
semantic feedback.

\begin{itemize}
\item C1: prompt-only. The natural-language rule is present, but there is no
ActPlane enforcement.
\item C2: audit. ActPlane observes and reports, but does not fail the action.
\item C3: fail without ActPlane reason. The violating action fails through
\texttt{-EPERM} or \texttt{SIGKILL}, but the agent receives no structured
remediation.
\item C4: fail with feedback. The action fails and ActPlane injects the rule
reason and remediation through the feedback file and hook bridge.
\end{itemize}

The main comparison is C3 to C4: the net benefit of the corrective feedback
loop. C1 to C2 estimates how often the prompt-only harness fails; C2 to C3
separates observation from action failure.

\subsection{Artifact Expectations}

The artifact should include policies, trigger scripts, agent prompts, raw logs,
statistical notebooks, and a container or VM description pinning kernel, BPF,
agent, and model versions. The final paper should avoid reporting any number
that cannot be regenerated from these artifacts.

\section{Related Work}

\subsection{In-Kernel Provenance and Information Flow}

ActPlane builds on a long line of provenance and information-flow systems.
PASS made provenance a storage-layer primitive~\cite{muniswamy2006pass}.
Hi-Fi and LPM used LSM hooks to capture trustworthy system
provenance~\cite{pohly2012hifi,bates2015lpm}. CamFlow captured whole-system
provenance streams over processes, files, and sockets~\cite{pasquier2017camflow}.

CamQuery is the closest prior system~\cite{pasquier2018camquery}. It executes
provenance queries inline in the kernel, propagates labels across process, file,
and network edges, and can prevent policy violations. ActPlane does not claim
that this mechanism is new. The difference is the domain and substrate:
ActPlane targets cooperative AI agents, compiles an agent-oriented source/sink
DSL to eBPF/BPF-LSM, and returns corrective feedback to the actor that violated
the rule.

TaintDroid, Dytan, libdft, and Panorama established dynamic taint tracking at
runtime, binary-instrumentation, or emulation layers~\cite{enck2010taintdroid,
clause2007dytan,kemerlis2012libdft,yin2007panorama}. ActPlane uses a much
coarser object-level taint. That coarseness is what allows kernel-wide coverage,
but it creates an over-taint risk that must be measured.

\subsection{eBPF and LSM Enforcement}

BPF-LSM, originally proposed through KRSI, lets eBPF programs attach to Linux
security hooks and return allow/deny verdicts~\cite{singh2019krsi,linuxbpflsm}.
ActPlane uses this as its pre-operation backend. Contemporary eBPF security
tools and proposals such as Tetragon, Tracee, and eBPF-PATROL demonstrate
practical event capture and enforcement, but generally evaluate per-event
predicates or lineage flags rather than typed cross-channel information
flow~\cite{ciliumtetragon,aquasecuritytracee,ghimire2025ebpfpatrol}.

\subsection{Agent Guardrails}

AgentSpec and Progent are the closest agent-policy systems at the tool
layer~\cite{wang2025agentspec,shi2025progent}. They provide deterministic
checks over tool names, arguments, or framework actions. ActPlane complements these systems by
moving the enforcement point below the tool API. If all effects pass through the
framework, tool-layer policies are sufficient and cheaper. If the agent can
execute arbitrary code, spawn subprocesses, or use direct system interfaces, the
OS layer is the robust boundary.

FIDES and CaMeL show that typed information-flow and capability separation are
useful abstractions for LLM agents~\cite{costa2025fides,debenedetti2025camel}. They operate
inside the agent loop or scaffolding. ActPlane applies the same style of
information-flow reasoning at the OS boundary.

\subsection{Agent Instruction Files}

Recent empirical work studies \texttt{CLAUDE.md}, \texttt{AGENTS.md}, and other
agent context files: their structure, content, maintenance, and effect on task
efficiency~\cite{chatlatanagulchai2025manifests,chatlatanagulchai2025agentreadmes,
santos2025claudeconfig,lulla2026agentsmd}. ActPlane's corpus study asks a
different question: which instructions are behavioral or flow invariants, and
which of those can be observed and enforced below the tool layer?

\section{Discussion}

\subsection{Security Tool or Harness?}

ActPlane can express security-relevant policies, such as keeping secret-derived
data local. However, its primary target is an agent harness: a deterministic
operating contract for a cooperative but forgetful agent. This affects design
choices. A classic reference monitor would prioritize pre-operation denial for
all security claims. ActPlane also needs audit-only guidance, harness-level
termination when pre-operation arguments are unavailable, and remediation text
that lets the agent complete the original task through an allowed path.

\subsection{Where the Claim Must Stay Honest}

The strongest novelty claim is not the taint algorithm. In-kernel label
propagation and provenance enforcement are established ideas. ActPlane must be
evaluated as a systems contribution that combines those ideas with a new
workload and interface: autonomous agents with rich tool access, natural-language
policy sources, and corrective feedback.

The final paper should be careful about three boundaries.

\textbf{Block is LSM-only.} Tracepoint observation cannot be described as
\texttt{block}. The current implementation correctly ignores \texttt{block}
rules in tracepoint mode and warns the operator.

\textbf{Prompt files are aspirational.} A corpus of instruction files shows what
developers want agents to do. It does not by itself prove how often agents
violate those rules. That is the job of the agent evaluation.

\textbf{Coarse taint can over-approximate.} Object-level labels are practical,
but they may carry more provenance than a byte-level system would. The evaluation
must measure false positives and demonstrate useful declassification and gate
patterns.

\section{Conclusion}

ActPlane moves agent harness enforcement below the tool layer. It compiles
agent-oriented information-flow policies into eBPF/BPF-LSM checks, tracks labels
across process, file, and endpoint objects, and returns kernel-detected
violations as corrective feedback to the agent. The evaluation plan is designed
to test bypass coverage, feedback-driven recovery, false positives, overhead,
and coverage over real agent instruction files without reporting numbers before
they can be reproduced.


\bibliographystyle{ACM-Reference-Format}
\bibliography{references}

\end{document}
